% GENERATED FILE. Edit canonical lessons, then run npm run generate:books.
% canonical-source-hash: fnv1a64:9daa6df8988e7d06
% canonical-lessons: TA-C153-eru, TA-C153-iranku, TA-C153-tirumpu, TA-C153-cer, TA-C153-kuli

\chapter{To climb, To get down, To turn, To join, To bathe}
\label{ch:ta-to-climb-to-get-down-to-turn-to-join-to-bathe}
\hlchaptermodality{153}

\begin{chapteropening}
\textbf{By the end of this chapter:} \emph{I can say to climb, to get down, to turn, to join and to bathe.}

Complete the last lesson of chapter 153: i can say to climb, to get down, to turn, to join and to bathe.
\end{chapteropening}

\section[ēṟu]{\ta{ஏறு} (ēṟu) — to climb, to get on}

\label{lesson:TA-C153-eru}



\begin{quote}
Before the new one: say the Tamil for to fly, then the Tamil for to fall.
\end{quote}

\subsection*{You'll want to know: \ta{ஏறு}}
\textbf{\ta{ஏறு}} — \emph{ēṟu} — \textquotedblleft{}to climb, to get on\textquotedblright{}.

\begin{grammarlens}[title={What you've built}]
One more word for this chapter.
\end{grammarlens}

\subsection*{Your turn}
\begin{itemize}
\raggedright
  \item \emph{Say it:} \emph{ēṟu}
  \item \emph{Say it:} \emph{ēṟu}, once more
  \item \emph{Say it:} say \emph{ēṟu}
  \item \emph{Recall:} say the Tamil for to fly, then the Tamil for to fall, then say \emph{ēṟu} again
\end{itemize}

\begin{tcolorbox}[breakable,colback=teal!4,colframe=teal!35!black,title={Before you move on}]
What does \ta{ஏறு} mean? (\textquotedblleft{}To climb, to get on\textquotedblright{}.) Say it once more.
\end{tcolorbox}
\section[iṟaṅku]{\ta{இறங்கு} (iṟaṅku) — to get down, to get off}

\label{lesson:TA-C153-iranku}



\begin{quote}
Before the new one: say the Tamil for to fall, then the Tamil for to climb.
\end{quote}

\subsection*{You'll want to know: \ta{இறங்கு}}
\textbf{\ta{இறங்கு}} — \emph{iṟaṅku} — \textquotedblleft{}to get down, to get off\textquotedblright{}.

\begin{grammarlens}[title={What you've built}]
One more word for this chapter.
\end{grammarlens}

\subsection*{Your turn}
\begin{itemize}
\raggedright
  \item \emph{Say it:} \emph{iṟaṅku}
  \item \emph{Say it:} \emph{iṟaṅku}, once more
  \item \emph{Say it:} say \emph{iṟaṅku}
  \item \emph{Recall:} say the Tamil for to fall, then the Tamil for to climb, then say \emph{iṟaṅku} again
\end{itemize}

\begin{tcolorbox}[breakable,colback=teal!4,colframe=teal!35!black,title={Before you move on}]
What does \ta{இறங்கு} mean? (\textquotedblleft{}To get down, to get off\textquotedblright{}.) Say it once more.
\end{tcolorbox}
\section[tirumpu]{\ta{திரும்பு} (tirumpu) — to turn, to return}

\label{lesson:TA-C153-tirumpu}



\begin{quote}
Before the new one: say the Tamil for to climb, then the Tamil for to get down.

Say \textbf{\ta{வா}}, then build on \emph{varu-}: \textbf{\ta{வருகிறேன்}}. Then \textbf{\ta{பார்}} and \textbf{\ta{கேள்}}.
\end{quote}

\subsection*{You'll want to know: \ta{திரும்பு}}
\textbf{\ta{திரும்பு}} — \emph{tirumpu} — \textquotedblleft{}to turn, to return\textquotedblright{}.

\begin{grammarlens}[title={What you've built}]
One more word for this chapter.
\end{grammarlens}

\subsection*{Your turn}
\begin{itemize}
\raggedright
  \item \emph{Say it:} \emph{tirumpu}
  \item \emph{Say it:} \emph{tirumpu}, once more
  \item \emph{Say it:} say \emph{tirumpu}
  \item \emph{Recall:} say the Tamil for to climb, then the Tamil for to get down, then say \emph{tirumpu} again
\end{itemize}

\begin{tcolorbox}[breakable,colback=teal!4,colframe=teal!35!black,title={Before you move on}]
What does \ta{திரும்பு} mean? (\textquotedblleft{}To turn, to return\textquotedblright{}.) Say it once more.
\end{tcolorbox}
\section[cēr]{\ta{சேர்} (cēr) — to join, to reach}

\label{lesson:TA-C153-cer}



\begin{quote}
Before the new one: say the Tamil for to get down, then the Tamil for to turn.
\end{quote}

\subsection*{You'll want to know: \ta{சேர்}}
\textbf{\ta{சேர்}} — \emph{cēr} — \textquotedblleft{}to join, to reach\textquotedblright{}.

\begin{grammarlens}[title={What you've built}]
One more word for this chapter.
\end{grammarlens}

\subsection*{Your turn}
\begin{itemize}
\raggedright
  \item \emph{Say it:} \emph{cēr}
  \item \emph{Say it:} \emph{cēr}, once more
  \item \emph{Say it:} say \emph{cēr}
  \item \emph{Recall:} say the Tamil for to get down, then the Tamil for to turn, then say \emph{cēr} again
\end{itemize}

\begin{tcolorbox}[breakable,colback=teal!4,colframe=teal!35!black,title={Before you move on}]
What does \ta{சேர்} mean? (\textquotedblleft{}To join, to reach\textquotedblright{}.) Say it once more.
\end{tcolorbox}
\section[kuḷi]{\ta{குளி} (kuḷi) — to bathe}

\label{lesson:TA-C153-kuli}



\begin{quote}
Before the new one: say the Tamil for to turn, then the Tamil for to join.
\end{quote}

\subsection*{You'll want to know: \ta{குளி}}
\textbf{\ta{குளி}} — \emph{kuḷi} — \textquotedblleft{}to bathe\textquotedblright{}.

\begin{grammarlens}[title={What you've built}]
One more word for this chapter.
\end{grammarlens}

\subsection*{Your turn}
\begin{itemize}
\raggedright
  \item \emph{Say it:} \emph{kuḷi}
  \item \emph{Say it:} \emph{kuḷi}, once more
  \item \emph{Say it:} say \emph{kuḷi}
  \item \emph{Recall:} say the Tamil for to turn, then the Tamil for to join, then say \emph{kuḷi} again
\end{itemize}

\begin{tcolorbox}[breakable,colback=teal!4,colframe=teal!35!black,title={Before you move on}]
What does \ta{குளி} mean? (\textquotedblleft{}To bathe\textquotedblright{}.) Say it once more.
\end{tcolorbox}
